\documentclass[11pt,twoside]{scrbook} % KOMA

% ---------- LINGUA & CODIFICA ----------
\usepackage[italian]{babel}
\usepackage[T1]{fontenc}

% ---------- TIPOGRAFIA ----------
\usepackage{lmodern}
\usepackage{microtype}

% ---------- GEOMETRIA PAGINA ----------
\usepackage[
  a4paper,
  top=24mm,
  bottom=28mm,
  inner=24mm,   % margine interno (verso rilegatura)
  outer=30mm,   % margine esterno: note a margine
  heightrounded
]{geometry}
\setlength{\headheight}{17pt} % evita sovrapposizioni con l'header
\setlength{\headsep}{6mm}     % distanza tra header e corpo

\usepackage[automark]{scrlayer-scrpage}
\clearpairofpagestyles
\automark[section]{chapter}
\ohead{\headmark}
\ofoot*{\pagemark}
\pagestyle{scrheadings}

% ---------- STRUMENTI DI STRUTTURA ----------
\setcounter{secnumdepth}{3}
\setcounter{tocdepth}{3}

% ---------- STILE TITOLI IN STILE KOMA ----------
\setkomafont{chapter}{\normalfont\huge\bfseries}
\setkomafont{section}{\normalfont\Large\bfseries}
\setkomafont{subsection}{\normalfont\large\bfseries}
\setkomafont{subsubsection}{\normalfont\normalsize\bfseries}

% ---------- STILE FRONTESPIZIO ----------
\setkomafont{title}{\sffamily\bfseries\Huge}
\setkomafont{author}{\sffamily\large}
\setkomafont{publishers}{\sffamily\normalsize}
\setkomafont{date}{\large}

\RedeclareSectionCommand[beforeskip=1.5ex,afterskip=1.5ex]{chapter}
\RedeclareSectionCommand[beforeskip=1ex,afterskip=0.8ex]{section}
\RedeclareSectionCommand[beforeskip=0.8ex,afterskip=0.6ex]{subsection}
\RedeclareSectionCommand[beforeskip=0.6ex,afterskip=0.4ex]{subsubsection}

% ---------- AMBIENTI UTILI ----------
\usepackage{amsmath,amssymb,amsthm}
\usepackage{booktabs}
\usepackage{array}
\usepackage{graphicx}
\usepackage{xcolor}
\usepackage{tcolorbox}
\usepackage{subcaption} 
\usepackage{upgreek} % \uptau
\usepackage{mathrsfs} % \mathscr{}
\usepackage{dsfont}
\usepackage{tikz}
\usetikzlibrary{intersections}
\usepackage{pgfplots} % grafici, creato in tikZ
\pgfplotsset{compat = 1.18}
\usetikzlibrary{
arrows.meta,
positioning,
calc,
decorations.markings,
angles,
quotes,
3d,
patterns
}
\usepackage[european]{circuitikz} % circuiti
\usepackage{comment} % \begin{comment}  \end{comment}
\usepackage{mathtools} % \coloneqq
\usepackage{siunitx}
\usepackage{hyperref}
\hypersetup{  % stile dei riferimenti
  colorlinks=true,
  linkcolor=black,
  citecolor=blue,
  urlcolor=blue
}
\sisetup{separate-uncertainty=true}
%----------- TEOREMI E STRUTTURE VARIE (KOMA, AMSTHM)

% Stile "teorema" (corsivo)
\theoremstyle{plain}
\newtheorem{theorem}{Teorema}[chapter]
\newtheorem{proposition}[theorem]{Proposizione}
\newtheorem{lemma}[theorem]{Lemma}
\newtheorem{corollary}[theorem]{Corollario}

% Stile "definizione" (testo dritto)
\theoremstyle{definition}
\newtheorem{definition}[theorem]{Definizione}

% Stile "nota" (più discorsivo)
\theoremstyle{remark}
\newtheorem{remark}[theorem]{Osservazione}
\newtheorem{example}[theorem]{Esempio}

% ---------- NOTE A MARGINE ----------
% \marginpar stampa nel margine esterno con oneside.
\newcommand{\mkeyword}[1]{%
  \marginpar{\raggedright\footnotesize\emph{#1}}%
}

% per le note a sinistra:
% \reversemarginpar

% ---------- PARAGRAFI E SPAZI ----------
\setlength{\parindent}{0.8em}
\setlength{\parskip}{0.25em}

%----------- COMANDI AGGIUNTIVI -----------

\newcommand{\1}{\mathds{1}} % in stile funzione caratteristica

% --- vettore colonna ---
\newcommand{\colvec}[1]{%
  \begin{pmatrix}
    #1
  \end{pmatrix}%
}

% --- derivata parziale ---
\newcommand{\der}[3][]{%
  \frac{\partial^{#1} {#2}}{\partial {#3}^{#1}}%
}

% --- versore --- 
\newcommand{\vers}[1]{
  \hat{\boldsymbol{#1}}
}

% --- epsilon ---
\newcommand{\eps}[0]{
  \varepsilon
}

% --- norma ---
\newcommand{\norm}[1]{
  \left\lVert #1 \right\rVert
}

% --- prodotto scalare ---
\newcommand{\sc}[1]{
  \langle #1 \rangle
}

% --- vettore ---
\renewcommand{\v}[1]{
  \mathbf{#1}
}

% --- \backslash ---
\newcommand{\bs}[0]{
	\backslash 
}

% --- \diag ---
\newcommand{\diag}[0]{
	\textnormal{diag} 
}

% --- Logaritmo principale ---
\newcommand{\Log}[0]{
  \textnormal{Log }
}

% --- overline ---
\newcommand{\ol}[1]{
  \overline{#1}
}

% --- underline ---
\newcommand{\ul}[1]{
  \underline{#1}
}

% --- valore assoluto ---
\newcommand{\abs}[1]{
  \left|{#1}\right|
}

% --- non andare a capo --- NON FUNZIONA
\newcommand{\inline}[1]{
  \binoppenalty=10000 \relpenalty=10000 #1
}

% --- Residui ---
\DeclareMathOperator{\Res}{
  \textnormal{Res}
}

% --- Supporto ---
\DeclareMathOperator{\supp}{
  \textnormal{supp}
}

% --- Overset ---
\newcommand{\ov}[2]{
  \overset{#2}{#1}
}

% --- Underset ---
\newcommand{\un}[2]{
  \underset{#2}{#1}
}

% --- Parte principale ---
\let\P\relax
\DeclareMathOperator{\P}{
  \mathcal{P}
}

% =================== Spazi ===================

% --- R numeri reali --- 
\newcommand{\R}[0]{
  \mathbb{R}
}

% --- C numeri complessi --- 
\newcommand{\C}[0]{
  \mathbb{C}
}

% --- E Funzioni lisce --- 
\newcommand{\E}[0]{
  \mathcal{E}
}

% --- S funzioni a decrescenza rapida --- 
\let\S\relax
\newcommand{\S}[0]{
  \mathcal{S}
}

% --- D Funzioni lisce a supporto compatto --- 
\newcommand{\D}[0]{
  \mathcal{D}
}

% =================== Trasformate ===================

% --- Fourier ---
\DeclareMathOperator{\F}{
  \mathcal{F}
}

% --- Laplace ---
\let\L\relax
\DeclareMathOperator{\L}{
  \mathcal{L}
}

% --- Hilbert ---
\let\H\relax
\DeclareMathOperator{\H}{
  \mathcal{H}
}

% =================== Scambi ===================

% --- phi φ ϕ ---
\let \ORIGINALphi \phi
\let \ORIGINALvarphi \varphi

\renewcommand{\phi}{\ORIGINALvarphi}
\renewcommand{\varphi}{\ORIGINALphi}

% =================== Stili ===================

% --- displaystyle ---
\newcommand{\DS}[0]{
	\displaystyle 
}

% --- textstyle ---
\newcommand{\TS}[0]{
	\textstyle 
}

% % --- scriptstyle --- nota: non attivi per mantere liberi i comandi \SS \SSS
% \renewcommand{\SS}[0]{
% 	\scriptstyle 
% }

% % --- \displaystyle ---
% \newcommand{\SSS}[0]{
% 	\scriptscriptstyle 
% }

% =================== Template ===================

% \mathrm{F}, \mathit{F}, \mathbf{F}, \mathsf{F}, \mathtt{F}, \boldsymbol{F}, \mathcal{F}, \mathscr{F}, \mathbb{F}, \mathds{F}, \mathfrak{F}

% ---------- BIBLIOGRAFIA ----------
\usepackage{csquotes} % \enquote{}
\usepackage[backend=biber,style=numeric,sorting=nyt]{biblatex}
\addbibresource{bibliografia.bib} % nome del file .bib

% ---------- FRONTESPIZIO ----------
\title{\textbf{\titolo}}
\author{Cosimo Pascarella}
\date{Anno Accademico 2026--2027}

\publishers{
  Email: \text{c.pascarella2@studenti.unipi.it}\\
  Versione: v0.1
}